\begin{frame}[noframenumbering,plain]
    \setcounter{framenumber}{1}
    \maketitle
\end{frame}

\begin{frame}
    \frametitle{Положения, выносимые на защиту}
    \begin{enumerate}
        \item Математическая модель задачи раскроя и~формирования производственного расписания для многомашинной производственной системы, включающая векторную целевую функцию, систему технологических, временных и~складских ограничений.
        \item Численный метод решения задачи на~основе генетического алгоритма с~событийной моделью формирования производственного расписания, обеспечивающий минимизацию отходов материала, числа раскроев и~переналадок оборудования.
        \item Комплекс программ для автоматизированного планирования раскроя и~формирования производственных расписаний в~целлюлозно-бумажной промышленности, реализующий разработанные модели и~численные методы.
    \end{enumerate}
\end{frame}
\note{
    Зачитываются положения, выносимые на защиту.
    Первое положение относится к математическому моделированию,
    второе~--- к численным методам,
    третье~--- к комплексу программ.
    Это соответствует трём составляющим специальности 1.2.2.
}

\begin{frame}
    \frametitle{Содержание}
    \tableofcontents
\end{frame}
\note{
    Работа состоит из четырёх глав.

    \medskip
    В первой главе проводится обзор литературы по задаче линейного раскроя и~смежным постановкам, описывается технологический контекст целлюлозно-бумажного производства и~формулируется постановка задачи.

    Во второй главе строится математическая модель: вводятся входные параметры, переменные решения (глобальный план), векторная целевая функция и~система из девяти ограничений.

    Третья глава посвящена численным методам: описан метод полного перебора (эталон), генетический алгоритм с~операторами мутации и~парного отбора, а~также алгоритм формирования расписания на~основе событийной модели.

    В четвёртой главе представлен комплекс программ, описаны форматы данных, интерфейсы и~результаты вычислительных экспериментов.
}
